\section{Why we did not train a policy}
\label{sec:nopolicy}

This section separates two things: the scientific conclusion the data supports, and
a project decision about resource commitment.

\subsection{Scientific conclusion}
\label{sec:scientific-conclusion}

\begin{enumerate}
  \item Multi-step planning improved over one-step control on held-out instances
        ($+0.456$~nat, CI $+0.254$ to $+0.720$, $35/40$).
  \item Replanning during execution did not meaningfully improve over committed
        planning: $+0.010$~nat, $21/40$.

        \begin{quote}
        The confidence interval was narrow and included zero, and we did not
        observe a practically large feedback benefit.
        \end{quote}

        We do not write ``equivalent'' or ``statistically the same''; there is no
        prior equivalence margin
        \citep{schuirmann1987tost,lakens2017equivalence}.
  \item Under stochastic model mismatch the committed plan was fragile. The
        rejection rate rose from $0.00$ to $0.66$--$0.79$ and terminal improvement
        fell below the tuned constant.
  \item Replanning reduced that fragility but did not consistently outperform
        inexpensive one-step control: $\text{replanning} - \text{one-step}$ was
        $-1.111$ and $-0.716$ in the two minibatch conditions
        ($n = 3$ per regime, exploratory).
\end{enumerate}

\subsection{Project decision}

\begin{quote}
PPO training was a conditional next stage, not a required component of the study. We
predeclared empirical gates to determine whether the added implementation and
computational cost was justified. Because feedback replanning did not consistently
outperform committed planning or inexpensive one-step control, we did not proceed to
PPO.
\end{quote}

That decision-search cost is $1{,}294$ times the deployment budget also enters this
judgement.

We state the distinction explicitly:

\begin{itemize}
  \item PPO did not fail.
  \item We did not run PPO.
  \item This is not a rejection of reinforcement learning in general.
  \item We did not establish headroom that would justify a learned state-dependent
        feedback policy in the present environment.
\end{itemize}

\subsection{The status of the gates}

The gate thresholds (for example $\gate{C3} \geq 0.3$~nat) are \emph{not} universal
theoretical criteria. Their character is as follows:

\begin{itemize}
  \item Policy learning requires a separate, large experimental stage.
  \item We therefore registered in advance a minimum effect size, a sign-consistency
        requirement, whether multi-step plans were used, and transfer conditions.
  \item The gates are a scope-control device, not a statistical truth verdict.
  \item We did not change the thresholds after seeing results.
\end{itemize}

The thresholds were fixed before the results were seen and were not changed
afterwards. Their particular values are a scope-control choice, not a claim of
correctness.

\subsection{Follow-up direction}
\label{sec:followup}

That \gate{C2} is GO on held-out instances means a good multi-step sequence exists,
and that \gate{C3} is small means that sequence can be determined at the initial
state. This combination points less towards a per-step feedback policy than towards
an \emph{amortized approach that predicts a schedule from initial problem features}
\citep{amos2023amortized,andrychowicz2016l2l,bae2022apo}. We did not implement it
and mention it only as a direction.
